\paragraph{Hardware and engine.}
All new GPU experiments use the same eight H200 replicas in Modal's us-west
region, with one GPU, eight CPU cores, and 64\,GiB of host memory per replica.
The base model is \texttt{Qwen/Qwen3.8-27B}, loaded as
\texttt{Qwen3\_5ForConditionalGeneration}; its hybrid attention/recurrent cache
is relevant to the matching procedure. Tensor, data, and pipeline parallelism
are each one. Context capacity is 16,384 tokens and the static memory fraction
is 0.8. The engine is SGLang v0.5.17 with the pinned Modal fork and routing
instrumentation (Appendix~\ref{sec:reproduce}). The cached model download records
revision \texttt{1d4bf0f2ff6012fd82039f2fa52739d0dd7c60c0}.
Its text configuration declares BF16; the engine's weight and KV dtype options
are \texttt{auto}. A live settings audit records FA3 attention, FlashInfer
sampling, CSGMV LoRA execution, CUDA graphs enabled, 8,192-token chunked prefill,
16,384 maximum prefill tokens, and page size one. The model configuration and
its SHA-256 are included in the artifact.

The engine uses longest-prefix-match scheduling, eight GPU adapter slots,
256 host adapter slots, and no overlapping LoRA loading. Adapters have rank 32
and update the query, key, value, output, gate, up, and down projections.
Although the configuration requests 128 running sequences, the deployed engine
reports an effective capacity of \textbf{36 per replica}, or 288 for the pool.
All capacity comparisons use that observed value. The Flash target concurrency
is 64 and the minimum and maximum replica counts are both eight. These are
fixed-pool experiments; they do not evaluate an autoscaling controller.

\paragraph{Calibration.}
A separate pre-experiment run measures aggregate rates on each replica and uses
their median: $r_p=8{,}706$ input tokens/s, $r_d=1{,}065$ output tokens/s,
$t_g=0.0672$\,s, and $t_c=1.871$\,s. These values remain fixed for the
comparisons. Registration measures exposed elapsed delay, including shared
storage and software overhead; it is not a guaranteed network-cold transfer.
The transfer estimate is the median backend-latency difference between paired
cold-GPU and hot-GPU requests with the same 128-token prompt and cache match.
Eight independent 8,192-token prompts estimate prefill rate; a warmed batch of
36 requests generating 512 tokens estimates decode rate. Separately, engine
load counters record host elapsed loading calls without CUDA synchronization.
Neither those summed counters nor the calibration differences are added
directly to a claimed wall-time saving.

These measurements also constrain what the score can plausibly optimize. One
GPU adapter load corresponds to approximately 585 uncached input tokens at the
measured prefill rate; one registration corresponds to approximately 16,287.
Reusing a 7,168-token prefix is therefore more valuable in this model than
avoiding one GPU load, but less valuable than avoiding one registration. These
are comparisons of calibrated work terms, not predictions of isolated request
latency. Repeated eviction and different batch shapes can change the realized
cost.

\paragraph{Checkpoints and clients.}
Sixteen independently trained DAPO math clients supply real checkpoint
identities and recorded update/publication intervals of approximately 34--96\,s.
Each delay is that source client's measured second update, from training start
to publication, reused as a constant in later replay rounds.
The RL-shaped replays use versions 4 and 7. A general-serving stress test uses
versions 4, 7, 10, and 13 from these jobs: 64 different checkpoint files,
verified by distinct SHA-256 digests. This represents 64 adapter versions, not
64 independent training jobs. Each checkpoint is approximately 319\,MB.

RL-shaped experiments use eight distinct CPU frontend containers (two CPU cores
and 8\,GiB each), two independent clients per frontend, and one router shared by
those two clients. A training client's requests originate from one frontend;
these GPU experiments do not test simultaneous same-adapter arrivals from many
frontends.
The dedicated benchmark frontends invoke the shared \texttt{sample\_task}
execution path directly. They bypass the public control-plane command queue;
reported request latency begins at frontend preparation, not at SDK submission.
Each client starts after a 0.5\,s incremental offset. After its initial rollout,
it overlaps the next rollout with its own recorded update delay. There are no
cross-client round barriers. At the next eligible round it switches to the
specified saved checkpoint. We validate independent progress from rollout
timestamps, not from launch staggering alone. Optimizers do not execute during
these replays; results concern inference under a training-shaped arrival process.

\paragraph{Controlled prompts and generations.}
\begin{table*}[t]\centering\small
\begin{tabular}{lrrrrrr}\toprule
Workload & Clients & Versions used & Sequence ceiling & Samples/group & Input/output tokens & Generations\\\midrule
Warmed reuse & 16 & 16 & 1,024 & 8 & 8,192 / 256 & 3,072\\
Version turnover & 16 & 32 & 1,088 & 8, 64, 128 & 8,192 / 256 & 4,352\\
Three-turn dialogue & 16 & 32 & 512 & $8\rightarrow1$ & 8,192--9,216 / 256 & 4,608\\
General serving & 16 & 64 & 512 & 1 & 8,192 / 512 & 2,048\\
Longer decode & 16 & $\le32$ & 1,088 & 8, 64, 128 & 8,192 / 512 or 1,024 & Variable\\
\bottomrule\end{tabular}
\caption{HTTP/2 workload definitions. ``Clients'' counts independent checkpoint
sources; general serving has eight streams per exact version. The sequence
ceiling is potential logical concurrency, not a claim that all streams are
simultaneously in an engine. All workloads share eight GPUs and 288 effective
running slots. Dialogue continuations use single-sequence requests. Longer
decode admits new rollouts for 300 seconds (512 outputs) or 180 seconds
(1,024 outputs), then drains already submitted work.}
\label{tab:workloads}\end{table*}
To control cost, long prompts contain 7,168 tokens shared by a client's problem
across rounds and a 1,024-token round-specific suffix. Tokens are sampled
deterministically from the vocabulary occurring in that client's DAPO prompts.
Thus the lengths and sharing structure are controlled, but the resulting long
text is synthetic. Output lengths are fixed with \texttt{ignore\_eos=True};
the principal RL comparisons generate 256 tokens per request. This is not an
evaluation of answer quality or the natural EOS-length distribution.

The warmed workload gives each of 16 clients eight problem groups and eight
samples per group: 1,024 potential concurrent sequences. The initial HTTP/1
phase uses four rounds (4,096 generations); the subsequent HTTP/2 comparisons
use three (3,072 generations), permitting repeated trials within the allocation.
Three serial one-token warmup requests per problem
establish reusable prefixes before timing. The version is fixed, while suffixes
change across rounds.

The turnover workload preserves the same total demand for 15 clients but
changes its concentration. Clients 1--7 each submit one group of 64 siblings;
client 0 submits one group of 128; clients 8--15 retain eight groups of eight.
The ceiling is 1,088 concurrent sequences and the four rounds contain 4,352
generations. Versions 4 and 7 serve two rounds each. Only version 4 is registered
initially; version 7 incurs real first-use registration during measurement.
This shape stresses a few popular affinity keys while other clients continue
independently.

The dialogue workload uses four groups of eight initial samples per client,
three dependent generation turns, and three rollout rounds. Its ceiling is
512 active sequences and it contains 4,608 generations. Each continuation
appends the actual preceding 256-token response and a 256-token synthetic tool
reply, producing input lengths 8,192, 8,704, and 9,216. The first group finishes
before its branches continue; later branches progress independently. Adapter
identity is fixed within a trajectory and changes between rollout rounds.

The final general-serving workload uses all 64 saved versions, eight independent
streams per version, four prefix groups per version, and four requests per
stream. It therefore has 512 concurrent streams and 2,048 requests. Inputs are
8,192 tokens with 7,168-token shared prefixes; outputs are 512 tokens. Each
stream waits for its previous request, but the returned text is not appended in
this workload. Eight independent router objects run in one eight-core CPU
container; this tests independent router state, not eight-process overhead.
Each router shard has its own HTTP/2 generation client, and a separate client
collects observations. These eight generation connections permit the intended
concurrency despite the client's per-connection stream cap.
All versions are registered before timing. The separate stress pilot doubles
the streams and uses three requests per stream (1,024 concurrency, 3,072 total).

RL sampling uses temperature 1 and deterministic per-request seeds derived from
4242, the round, group, and sibling index. General serving uses temperature 0
and sampling seed $12026+100\times\text{stream index}+\text{turn}$.
Prompts, checkpoint identities, and requested output limits are
identical across policies within each trial.

The longer-decode controls retain the version-turnover shape and run independent
loops for a 300\,s rollout-admission window with 512-token outputs and a 180\,s
window with 1,024-token outputs. The latter window is shortened before launch
to fit the unchanged GPU deadline, equally for both policies.
Already submitted rollouts finish after that window. We report completions in
the interval from 60\,s to each window's end separately from total replay time
and drain time.
Because faster clients start subsequent rollouts earlier, these are closed-loop
rate measurements with policy-dependent request totals, not fixed-work makespan
comparisons. A 100-round guard is not reached in the measured window.

\paragraph{Comparisons and reset.}
The main controls are native KV affinity (KV), KV with cohort preparation
(KV+P), and stable adapter domains with preparation (Domains+P). The turnover
experiment also includes unrestricted estimated-work routing with preparation
(Work+P). General serving additionally tests the default SDK path with no
caller-supplied cache-affinity key; its ordinary per-request session header
remains present. Thus ``default'' does not mean removing every routing header.
The initial comparison orders KV before KV+P before Domains+P;
repeated comparisons reverse the latter pair. Every final trial unregisters the
entire previous adapter catalog, registers the prescribed initial versions,
materializes the base-model health-probe slot, and flushes prefix state.
Warmup, registration, and frontend startup are outside the measurement interval.
On-demand registration of new versions and cohort preparation are inside it.

\paragraph{Transport amendment.}
The initial stress trials use HTTP/1 with a 240\,s timeout; the subsequent
four-policy turnover and warmed trials increase it to 600\,s. Warmed trials
still encounter roughly 360\,s read failures after requests disappear from
engine observations. We preserve these runs and stop the unfinished HTTP/1
suite. A separate phase enables HTTP/2 for every arm, validates greedy cohort
equivalence again, and asserts that every successful physical generation
attempt records HTTP/2 trace events. The server advertises 1,024 concurrent
streams. We do not compare throughput across transport phases as a routing
effect, and do not claim that HTTP/2 alone establishes the root cause of the
HTTP/1 failures. The allocation deadline is extended by 45 minutes to run this
controlled phase; the GPU count remains eight.

The first general-serving HTTP/2 attempt exposed an additional client-side
limit. HTTPcore 1.0.9 permits only 100 streams on one connection, despite the
server's higher limit. That attempt shared one connection between generation
and observation, reached only 82 observed active requests, and had observation
gaps up to 180\,s. We preserve its completed default arm as a diagnostic and
restart every general-serving policy with eight generation connections and a
separate observer. The RL replays already have eight independent frontend
clients. Corrected general-serving arms must observe at least 288 active pool
admissions; this validates pressure at the pool boundary, not GPU utilization.
The RL frontends retain the same 100-stream-per-connection limit under every
policy. Their logical concurrency ceilings can therefore exceed actual HTTP
concurrency. In the first trials, measured peak pool admissions are 461--480
for warmed reuse, 542--600 for turnover, and 375--400 for dialogue, all above
the pool's effective running capacity. We report both configured demand and
observed pressure rather than equating them.

\paragraph{Metrics and validation.}
Useful throughput counts only requested output tokens. Request latency begins
before cohort preparation and ends when that individual sample completes;
rollout latency ends after the client's entire batch or dialogue finishes.
Replay makespan includes the recorded update intervals, including the final
overlapped interval. We therefore report rollout latency alongside replay
throughput rather than interpreting all elapsed-time changes as engine capacity.

KV hit ratio is reused prompt tokens divided by requested prompt tokens.
We also account separately for auxiliary seed prefill, which can raise the
visible hit ratio without eliminating the initial computation. Hot-LoRA ratio
is the fraction of engine admissions finding the adapter GPU-resident; it
includes auxiliary seeds and physical retry admissions. Adapter load and
eviction counts come from engine counters, not router predictions. Replica load
is measured at two-second intervals; busy-balance statistics require at least
288 active pool requests. They are sampled observations, not continuous traces.

Every completed arm verifies exact start/end adapter versions, output and
log-probability lengths, finite log probabilities, unique request identities,
destination consistency, and zero final sidecar inputs. Native final queue
counts are zero except where a cached subsecond observation predates that
replica's last successful response; such cases require no retries and are
reported without changing the raw counts. Later arms poll for settled native
zeros after freezing the timer. The cohort control separately
checks that greedy outputs agree before and after preparation on all eight
replicas. Greedy token equality does not establish identical floating-point
log probabilities or training-quality equivalence. Retries, incomplete arms,
observation gaps, and exploratory trials are retained and identified explicitly.
